% =====================================================================
%  MATH300 Project Description -- Algorithms Behind the Libraries
% =====================================================================
\documentclass[11pt,a4paper]{article}

\usepackage[T1]{fontenc}
\usepackage[utf8]{inputenc}
\usepackage{lmodern}
\usepackage[margin=2.5cm]{geometry}
\usepackage{enumitem}
\usepackage{booktabs}
\usepackage{tabularx}
\usepackage{array}
\usepackage{amsmath,amssymb}
\usepackage{xcolor}
\usepackage[colorlinks=true,urlcolor=blue,linkcolor=black,citecolor=black]{hyperref}

\setlist{itemsep=2pt, topsep=4pt}
\setlength{\parindent}{0pt}
\setlength{\parskip}{6pt}

% Placeholder macro: blue text in square brackets.
% To hide the colour in the final version, change {blue} to {black}.
\newcommand{\ph}[1]{\textcolor{blue}{[#1]}}

% Marker for optional bonus tasks.
\newcommand{\bonus}{\textbf{(Bonus)}}

\newcolumntype{L}{>{\raggedright\arraybackslash}X}

\begin{document}

% ---------------------------------------------------------------------
%  Title (no logo, no instructor name)
% ---------------------------------------------------------------------
\begin{center}
  {\large\bfseries MATH300: Mathematical Computing \textperiodcentered{} Fall 2026}\\[6pt]
  {\Large\bfseries Project: Inside the Black Box, Algorithms Behind the Libraries}
\end{center}
\vspace{6pt}

Groups must be formed, and project titles submitted on Canvas, by
\textbf{October 26, 2026}. The project
has two graded deliverables and an oral exam: a 3--5 page report for
Project~2 (due \textbf{November 20, 2026}), a poster for the final project
(due \textbf{December 11, 2026}), and an oral exam held between November 30
and December 17, 2026, 3:30\,pm. Read the general rules below before you
start; they apply throughout the project.

% ---------------------------------------------------------------------
\section*{General remarks}
% ---------------------------------------------------------------------
\begin{itemize}
  \item Projects are set for a group of \textbf{two or three} students.
  \item Form your group by \textbf{October 26, 2026}, and submit your
        project title as a group assignment on \textbf{Canvas} by the same
        date.
  \item Every group member must carry an \textbf{equal portion} of the coding
        and of the report and poster preparation.
  \item The scope can be adjusted to the number of participants, their
        interests and prior knowledge.
  \item Only \textbf{Python} may be used for programming.
  \item Implementation, documentation, validation and demonstration examples
        are all part of the project.
  \item Within these rules you are free to choose the implementation, code
        structure and exact content. You can discuss this with me at any time.
\end{itemize}

% ---------------------------------------------------------------------
\section*{Version control and GitHub}
% ---------------------------------------------------------------------
Each group keeps \textbf{one GitHub repository for the whole project}
(Project~2 and the final project), and every member commits from their own
GitHub account.
\begin{itemize}
  \item Use git from the very beginning. Commits should be small, consistent
        and clearly described. Avoid large commits.
  \item Each group member uploads their own code using \textbf{their own
        GitHub account}. Commits from a shared or borrowed account will not
        be credited to the other members.
  \item The repository name is up to your group.
  \item If the repository is private, add my GitHub account
        \textbf{\texttt{say-yas}} as a collaborator.
  \item The repository must contain a README describing the project, how to
        install and run the code, and the AI-use statement (see below).
\end{itemize}

% ---------------------------------------------------------------------
\section*{Use of AI tools}
% ---------------------------------------------------------------------
\textbf{AI-generated code is not accepted in any form.} Code produced by a
chat assistant or code generator (e.g.\ ChatGPT, Claude, Gemini) and copied,
pasted or adapted into your project will be treated as work that is not your
own.
\begin{itemize}
  \item \textbf{Permitted:} inline coding assistance such as GitHub Copilot
        autocompletion while you write the code yourself, \textbf{provided its
        use is acknowledged}.
  \item \textbf{Acknowledgement:} include an AI-use statement in the README
        and in the report that names the tool(s) used and the parts of the
        code (files or functions) where they assisted. If no AI tools were
        used, state this explicitly.
  \item \textbf{Reports and posters must not be generated by AI} in any
        form: their text, figures and layout must be your own work. If AI
        tools were used for any other non-code purpose (e.g.\ checking
        grammar or searching the literature), describe this in the same
        statement.
  \item You are responsible for every line in your repository. In the oral
        exam you must be able to explain and modify any part of the code you
        submitted.
\end{itemize}

% ---------------------------------------------------------------------
\section*{Deliverables and deadlines}
% ---------------------------------------------------------------------
Submit each deliverable following the details provided on the
\textbf{Canvas} page. The written report and poster are handed in on
\textbf{Crowdmark}, and the link to your GitHub repository is submitted
separately on \textbf{Canvas}.

\begin{center}
\small
\renewcommand{\arraystretch}{1.35}
\begin{tabularx}{\textwidth}{|>{\raggedright\arraybackslash}p{2.5cm}|>{\raggedright\arraybackslash}p{3.0cm}|>{\raggedright\arraybackslash}p{3.0cm}|L|}
\hline
\textbf{Deliverable} & \textbf{Due} & \textbf{Submission} & \textbf{Notes} \\
\hline\hline
Group and project title & Monday, October 26, 2026 &
  Group assignment on Canvas &
  Group of two or three members; submit the project title. \\
\hline
Project~2 report & Friday, November 20, 2026 &
  Report (PDF) on Crowdmark; repository link on Canvas &
  3--5 pages; the page limit \textbf{includes} references, figures and
  appendices. \\
\hline
Final project poster & Friday, December 11, 2026 &
  Poster (PDF) on Crowdmark; repository link on Canvas &
  36 inches high by 48 inches wide (landscape). Target audience:
  undergraduate classmates with a basic background in math, physics or
  engineering. \\
\hline
Oral exam & Between November 30 and December 17, 2026, 3:30\,pm &
  Individual meeting; details discussed in class &
  30 minutes per student. Your poster must already be submitted on
  Crowdmark before your exam. The last possible slot ends December 17 at
  3:30\,pm. \\
\hline
\end{tabularx}
\end{center}

\subsection*{Oral exam}
The oral exam is individual and lasts 30 minutes per student. It takes place
between November 30 and December 17, 2026, 3:30\,pm. The oral exams and how
to arrange them will be discussed in class before the Thanksgiving break.

Before your oral exam, your poster must also be submitted on
\textbf{Crowdmark}; that is, you must have finished both Project~2 and the
final project. In the meeting you will:
\begin{itemize}
  \item Present your results and algorithms.
  \item Possibly live-code and run parts of your algorithms to generate
        outputs on the spot.
  \item Answer questions about the theory, your implementation and your
        results. Each member should be able to explain the whole project, not
        only their own parts.
\end{itemize}
Bring a laptop with the repository cloned and a working environment, so that
the code runs without additional setup.

% ---------------------------------------------------------------------
\section*{Grading}
% ---------------------------------------------------------------------
\begin{center}
\small
\renewcommand{\arraystretch}{1.35}
\begin{tabularx}{0.8\textwidth}{|L|>{\raggedleft\arraybackslash}p{3.2cm}|}
\hline
\textbf{Component} & \textbf{Points} \\
\hline\hline
Group and project title (on time) & 10 \\
\hline
Project~2 report & 100 \\
\hline
Final project poster & 100 \\
\hline
Oral exam (individual) & 50 \\
\hline
Bonus: extra steps beyond the required tasks & up to 10 \\
\hline
\end{tabularx}
\end{center}

The Project~2 report and the poster are each graded out of 100. What matters
most is:
\begin{itemize}
  \item \textbf{Coherence:} a clear story from the question, through the
        methods, to the results and conclusions.
  \item \textbf{Correctness:} of the mathematics, the code and the results.
  \item \textbf{Depth of analysis and discussion:} explain what your results
        mean, why they look the way they do, and what their limitations are,
        rather than only showing them.
  \item \textbf{Professional presentation:} figures with labeled axes and
        readable font sizes, captions and legends, a clean and readable
        layout, and proper references.
\end{itemize}

A complete, correct and well-discussed
solution of the required tasks can earn full marks. Tasks marked
\bonus{} and the suggestions under ``Going further'' are optional; exploring
such extra steps can earn \textbf{up to 10 bonus points} in total.

% ---------------------------------------------------------------------
\section*{Requirements for the report and the poster}
% ---------------------------------------------------------------------
Both the report and the poster must contain a clear description of the
project, references, properly labeled figures and an author-contributions
statement.

\subsection*{Content}
\begin{itemize}
  \item A clear and concise description of the problem, the methods and the
        results.
  \item The AI-use statement described above.
  \item The report should not contain any significant amount of code; refer
        to the repository instead.
\end{itemize}

\subsection*{Poster}
\begin{itemize}
  \item The poster must be \textbf{self-contained}: a reader should be able to
        follow it without the report, the code or a spoken explanation.
  \item You may be selective in what you show and how you narrate your story,
        so that an undergraduate audience with a basic background in math,
        physics or engineering can follow it.
  \item You may choose your own title for the poster; be creative.
\end{itemize}

\subsection*{References}
\begin{itemize}
  \item Include a list of references at the end of the report and the poster.
  \item Cite sources clearly where they are used in the text (papers, books,
        software, code repositories and data sets).
  \item Use one consistent citation style throughout.
\end{itemize}

\subsection*{Figures}
\begin{itemize}
  \item Every plot must have labeled $x$ and $y$ axes (and $z$ axis for 3D
        plots or color bars), including units where applicable.
  \item Axis labels, tick labels and legends must have a font size similar to
        the surrounding text, so they remain readable without zooming.
  \item Each figure needs a legend when it shows more than one data series,
        and a caption that states what is shown.
  \item Use log scales where appropriate (e.g.\ for convergence plots) and say
        so in the caption.
\end{itemize}

\subsection*{Author contributions}
Each group must include an author-contributions statement in the report and
the README that states:
\begin{itemize}
  \item Who did what: which parts of the code, experiments, analysis, writing
        and poster each member worked on.
  \item To what extent: the approximate share of each member in each part
        (e.g.\ 50\,\% / 50\,\%).
  \item Who took the lead in which parts.
\end{itemize}
The coding and the report and poster preparation must be shared equally
between the members. The commit history should be consistent with the stated
contributions.

% ---------------------------------------------------------------------
\section*{Codebase requirements}
% ---------------------------------------------------------------------
\begin{itemize}
  \item Well structured and documented, both in the code (docstrings,
        comments) and in the repository (README).
  \item Validated: show that your implementation produces correct results,
        e.g.\ on examples with known solutions.\end{itemize}

% ---------------------------------------------------------------------
\section*{Project description}
% ---------------------------------------------------------------------

\subsection*{Goal}
In class we have used library functions to interpolate data, find roots and
decompose matrices, usually without looking at what happens inside them. In
this project you will open these black boxes. You will choose \textbf{one theme}
from the list below (for example interpolation) and \textbf{at least two
methods} within that theme, explain how each one works, present its pseudocode, implement it yourself in Python, and
compare your implementation with the version available in a package such as
NumPy or SciPy \cite{scipy}.

\subsection*{Possible themes and methods}
Choose \textbf{one theme} from the following list and \textbf{at least two
methods} within it. You may also propose other methods or another theme
related to the course (please check them with me first).
\begin{itemize}
  \item \textbf{Interpolation:} cubic spline interpolation \cite{deboor};
        PCHIP, the monotone piecewise cubic Hermite interpolation of Fritsch
        and Carlson \cite{fritsch}; Akima interpolation \cite{akima}. SciPy:
        \texttt{CubicSpline}, \texttt{PchipInterpolator},
        \texttt{Akima1DInterpolator} \cite{scipyinterp}.
  \item \textbf{Root finding:} Ridders' method \cite{ridders}; Brent's method
        \cite{brent}; the TOMS~748 algorithm of Alefeld, Potra and Shi
        \cite{toms748}. SciPy: \texttt{ridder}, \texttt{brentq},
        \texttt{toms748} \cite{scipyopt}.
  \item \textbf{Matrix decompositions:} QR decomposition by (modified)
        Gram--Schmidt or Householder reflections; LU decomposition with
        partial pivoting; Cholesky decomposition; SVD, for example by the
        one-sided Jacobi method \cite{trefethen,golub}. NumPy and SciPy:
        \texttt{numpy.linalg.qr}, \texttt{scipy.linalg.lu},
        \texttt{numpy.linalg.cholesky}, \texttt{numpy.linalg.svd}
        \cite{numpylinalg}.
\end{itemize}

\subsection*{Background}
All three root-finding methods above start from an interval $[a, b]$ on which
$f$ changes sign, so that it contains a root, and shrink this interval step
by step. The simplest such method, bisection, halves the interval in every
step; the methods above use cleverer guesses (for example by fitting a
curve through known points) to converge much faster, while still keeping the
root safely inside the interval.

Interpolation constructs a function that passes exactly through given data
points $(x_i, y_i)$. The three interpolation methods above all use a cubic
polynomial between each pair of neighboring points; they differ in how they
choose the slopes at the data points. A cubic spline makes the whole curve
twice continuously differentiable, which gives a very smooth curve but can
overshoot. PCHIP chooses the slopes so that the curve does not overshoot
monotone data, and Akima's method chooses them from nearby points only, which
reduces wiggles near outliers.

A matrix decomposition writes a matrix as a product of simpler matrices,
for example $A = QR$ with $Q$ orthogonal and $R$ upper triangular, $A = LU$
with $L$ lower and $U$ upper triangular, or $A = U S V^{\mathsf{T}}$ (the
SVD). These decompositions are the workhorses behind solving linear systems,
least-squares fitting and data compression. Different algorithms for the
same decomposition can behave very differently in floating-point arithmetic:
for example, classical Gram--Schmidt can lose the orthogonality of $Q$,
while Householder reflections keep it.

\subsection*{Tasks for Project 2 (report, due November 20, 2026)}
For \textbf{each} of your (at least two) chosen methods:
\begin{enumerate}
  \item Explain in your own words how the method works, and cite your
        sources.
  \item Present its pseudocode. (Pseudocode is welcome in the report;
        Python code is not.)
  \item Implement it yourself in Python. Do not call the library version
        inside your own implementation.
  \item Check your implementation on a few examples with known answers, and
        compare its results with the library version.
  \item \bonus\ Construct an example where one of your methods fails or
        struggles, and explain why.
  \item \bonus\ Implement a third method from the same theme, and include it
        in the comparison in the final project.
\end{enumerate}

\subsection*{Tasks for the final project (poster, due December 11, 2026)}
\begin{enumerate}
  \item Compare your methods with each other and with the library versions
        on several test problems: accuracy, and number of iterations or
        function evaluations where this applies. Present the results in
        plots or tables.
  \item Discuss what you observe: when does each method work well, and when
        does it struggle? Which method would you recommend for which kind of
        problem? Why might the library version be faster or more robust than
        yours?
  \item \bonus\ Compare the computing time of your implementations with the
        library versions.
\end{enumerate}

\subsection*{Going further}
These optional steps count toward the bonus points. For example, you can
implement a method from another theme.

% ---------------------------------------------------------------------
\begin{thebibliography}{9}
  \bibitem{scipy} SciPy documentation. \url{https://docs.scipy.org/doc/scipy/}
  \bibitem{scipyinterp} SciPy. \emph{Interpolation (scipy.interpolate)}.
    \url{https://docs.scipy.org/doc/scipy/reference/interpolate.html}
  \bibitem{scipyopt} SciPy. \emph{Optimization and root finding
    (scipy.optimize)}.
    \url{https://docs.scipy.org/doc/scipy/reference/optimize.html}
  \bibitem{deboor} C.~de~Boor (2001). \emph{A Practical Guide to Splines},
    revised edition. Springer.
    \url{https://link.springer.com/book/9780387953663}
  \bibitem{fritsch} F.~N.~Fritsch and R.~E.~Carlson (1980). Monotone
    piecewise cubic interpolation. \emph{SIAM Journal on Numerical
    Analysis}, 17, 238--246. \url{https://doi.org/10.1137/0717021}
  \bibitem{akima} H.~Akima (1970). A new method of interpolation and smooth
    curve fitting based on local procedures. \emph{Journal of the ACM}, 17,
    589--602. \url{https://doi.org/10.1145/321607.321609}
  \bibitem{ridders} C.~J.~F.~Ridders (1979). A new algorithm for computing a
    single root of a real continuous function. \emph{IEEE Transactions on
    Circuits and Systems}, 26, 979--980.
    \url{https://doi.org/10.1109/TCS.1979.1084580}
  \bibitem{brent} R.~P.~Brent (1973). \emph{Algorithms for Minimization
    without Derivatives}. Prentice-Hall.
    \url{https://maths-people.anu.edu.au/~brent/pub/pub011.html}
  \bibitem{toms748} G.~E.~Alefeld, F.~A.~Potra and Y.~Shi (1995). Algorithm
    748: Enclosing zeros of continuous functions. \emph{ACM Transactions on
    Mathematical Software}, 21, 327--344.
    \url{https://na.math.kit.edu/alefeld/download/1995_Algorithm_748_Enclosing_Zeros_of_Continuous_Functions.pdf}
  \bibitem{trefethen} L.~N.~Trefethen and D.~Bau (1997). \emph{Numerical
    Linear Algebra}. SIAM. \url{https://doi.org/10.1137/1.9780898719574}
  \bibitem{golub} G.~H.~Golub and C.~F.~Van~Loan (2013). \emph{Matrix
    Computations}, 4th edition. Johns Hopkins University Press.
    \url{https://nhigham.com/2013/05/31/fourth-edition-of-matrix-computations/}
  \bibitem{numpylinalg} NumPy. \emph{Linear algebra (numpy.linalg)}.
    \url{https://numpy.org/doc/stable/reference/routines.linalg.html}
\end{thebibliography}

\end{document}
